\section{Open problems}\label{sec:open}

Theorem A answers \cref{prob:q91} for linear codes over a fixed field of characteristic $127$ and for binary codes, and
Corollary U answers its undirected form. The following questions remain, roughly in order of their importance for
linear-time transposition.

\begin{enumerate}[label=(\arabic*)]
\item \emph{Native relations.} Does the half bound, or a fixed deficit below rate $1$, hold for the actual cone relations
of oblivious linear-time machines, as in the conclusion of~\cite[Proposition~9.2]{DirectedRecovery2026}? Which property
of tape computations, beyond bounded degree and paired distances, excludes twisted separated structures, or does some
machine realize them?
\item \emph{A deficit.} Is there, for every $\Delta$, a $\rho(\Delta)>0$ such that storage codes on relations of incidence
degree at most $\Delta$ with the paired-distance property have rate at most $1-\rho(\Delta)+o(1)$? The rates that our bounds
guarantee approach $4-2\sqrt3\approx0.536$; codes of rate tending to $1$ on such relations would answer this negatively.
\item \emph{Characteristic $2$ and $\F_2$-linear codes.} Is there a counterexample with $\F_2$-linear codes? The prime $2$ is
inert in $\Q(\sqrt5)$, but the tree then has degree $5$, far too small for \cref{thm:D}; a construction would need trees
of large degree $2^s+1$ with a collapse of the same tensor type.
\item \emph{Small degree.} What is the least $\Delta$ at which \cref{prob:q91} has a negative answer, and the least
degree at which the fourth paper's Question~6.1 has one? Our bounds $\Delta<127^{8\,791\,448}$ and $2\Delta$ come from the depth
of the order in \cref{thm:D}. For recovery sets of size at most three the seventh paper's Conjecture~P would give an
affirmative answer to \cref{prob:q91}~\cite{NestedElimination2026}, and for linear codes on triangle-free graphs of
maximum degree three the fourth paper proves the half bound~\cite[Theorem~6.13]{CurvedMotion2026}.
\item \emph{Beyond $4-2\sqrt3$.} Do higher-dimensional Goren--Oort configurations, or other algebraic sources of
tensor-type collapse, give rates above $4-2\sqrt3$ on bounded-degree relations with the paired-distance property?
\item \emph{Characteristic zero.} Do linear codes over fields of characteristic zero obey the half bound on all
bounded-degree relations with the paired-distance property, without the coefficient and width bounds of the fifth paper's
complex-linear theorem (\cref{rem:char0})?
\end{enumerate}
